\section{A detailed proof of perverse exactness}
\label{sec:fourier-perverse-exactness-detail}

We now expand the perverse-exactness capsule
\cref{found:fourier-perverse}.  The only input beyond the six operations and
perverse Artin vanishing is the following support theorem.  It is the step
that permits Verdier duality to be used without changing the Fourier
transform from a \(!\)-pushforward to a \(*\)-pushforward.

\begin{theorem}[Laumon's support-forgetting theorem]
\label{thm:fourier-support-forgetting}
For \(\cK\in\Dbc(A_x,\Qbarl)\), the canonical morphism
\[
 Rp_{y!}\bigl(p_x^*\cK\otimes\mu^*\cL_\psi\bigr)
 \longrightarrow
 Rp_{y*}\bigl(p_x^*\cK\otimes\mu^*\cL_\psi\bigr)
\tag{10.15}\label{eq:fourier-support-forgetting}
\]
is an isomorphism.
\end{theorem}

\begin{proof}
Compactify the \(x\)-line by
\(a:\A^1_x\hookrightarrow\mathbf P^1_x\), let
\(\eta_\infty\) be the generic point of the strict henselian trait at
\(\infty_x\), and let \(I_\infty\) be its inertia group.  The localization
triangle for this compactification identifies the cone of
\eqref{eq:fourier-support-forgetting} canonically with the
inertia-cohomology contribution
\[
 R\Gamma\!\left(
 I_\infty,\,
 \cK_{\bar\eta_\infty}\otimes
 R\Psi_{\infty_x}\bigl((a\times1)_!\mu^*\cL_\psi\bigr)
 \right).
\tag{10.17}\label{eq:fourier-boundary-inertia}
\]
Here the nearby-cycle factor is a complex over the \(y\)-line; this is a
global family statement, not a comparison of \(!\)- and \(*\)-cohomology on
each individual fiber.

Laumon's universal local-acyclicity theorem gives
\[
 R\Psi_{\infty_x}\bigl((a\times1)_!\mu^*\cL_\psi\bigr)=0.
\]
Consequently the complex in
\eqref{eq:fourier-boundary-inertia} vanishes, so the support-forgetting cone
vanishes.  This is \cite[Thm.~1.3.1.1]{LaumonFourier}; the displayed
nearby-cycle vanishing is Theorem 1.3.1.2 and Remark 1.3.1.3 there.
\end{proof}

\begin{proposition}[Fourier transform and Verdier duality]
\label{prop:fourier-verdier-detailed}
There is a functorial isomorphism
\[
 \D_{A_y}\FT_\psi(\cK)
 \simeq
 \FT_{\psi^{-1}}\bigl(\D_{A_x}\cK\bigr)(1).
\tag{10.16}\label{eq:fourier-verdier-detailed}
\]
\end{proposition}

\begin{proof}
Put
\(\cM=p_x^*\cK\otimes\mu^*\cL_\psi\).
Relative Verdier duality and the shift in the definition of \(\FT_\psi\)
give
\[
 \D_{A_y}\bigl(Rp_{y!}\cM[1]\bigr)
 \simeq Rp_{y*}\D_{A_x\times A_y}(\cM)[-1].
\]
The projection \(p_x\) is smooth of relative dimension \(1\), and
\(\cL_\psi^\vee=\cL_{\psi^{-1}}\).  Thus
\[
 \D_{A_x\times A_y}(\cM)
 \simeq
 p_x^*(\D_{A_x}\cK)\otimes\mu^*\cL_{\psi^{-1}}[2](1).
\]
Substitution yields
\[
 Rp_{y*}\bigl(
 p_x^*(\D\cK)\otimes\mu^*\cL_{\psi^{-1}}\bigr)[1](1).
\]
Apply \cref{thm:fourier-support-forgetting}, with
\(\psi^{-1}\) in place of \(\psi\), to replace \(Rp_{y*}\) by
\(Rp_{y!}\).  The result is exactly
\eqref{eq:fourier-verdier-detailed}.  This is Laumon
1.3.2.2; the twist is forced by
\(R\Gamma_c(\A^1,\Qbarl)=\Qbarl(-1)[-2]\), the same normalization check used
in Fourier inversion.
\end{proof}

\begin{theorem}[Perverse \(t\)-exactness of Fourier--Deligne transform]
\label{thm:fourier-perverse-texact-detailed}
The functor
\[
 \FT_\psi:\Perv(A_x,\Qbarl)\longrightarrow\Perv(A_y,\Qbarl)
\]
is exact.  Together with Fourier inversion, it is an equivalence of abelian
categories with quasi-inverse \(\FT_{\psi^{-1}}(1)\).
\end{theorem}

\begin{proof}
Let \(\cP\) be perverse on \(A_x\).  Since \(p_x\) is smooth of relative
dimension \(1\), smooth perverse pullback gives
\[
 p_x^*\cP[1]\in\Perv(A_x\times A_y).
\]
Tensoring with the lisse rank-one sheaf
\(\mu^*\cL_\psi\) preserves perversity.  The projection \(p_y\) is affine,
so perverse Artin vanishing says that \(Rp_{y!}\) is left \(t\)-exact.
Using the defining shift inside the displayed perverse object, we obtain
\[
 \FT_\psi(\cP)
 =Rp_{y!}\bigl(p_x^*\cP[1]\otimes\mu^*\cL_\psi\bigr)
 \in{}^p\D^{\geq0}(A_y).
\]

Verdier duality is \(t\)-exact for the middle perverse structure, so
\(\D\cP\) is perverse.  Applying the preceding lower bound with
\(\psi^{-1}\), and then \cref{prop:fourier-verdier-detailed}, gives
\[
 \D\FT_\psi(\cP)
 \simeq\FT_{\psi^{-1}}(\D\cP)(1)
 \in{}^p\D^{\geq0}(A_y).
\]
Because duality exchanges
\({}^p\D^{\leq0}\) and \({}^p\D^{\geq0}\), this is equivalent to
\(\FT_\psi(\cP)\in{}^p\D^{\leq0}(A_y)\).  Hence the transform lies in the
perverse heart.  A triangulated \(t\)-exact functor induces an exact functor
on hearts.  Finally \cref{thm:fourier-inversion} identifies
\(\FT_{\psi^{-1}}(1)\) as its quasi-inverse.  This is
\cite[Thm.~1.3.2.3]{LaumonFourier}.
\end{proof}

\begin{perversebox}
The argument uses both one-sided affine bounds.  The \(!\)-definition gives
the lower perverse bound directly; the upper bound is obtained only after
duality and support-forgetting.  Treating \(Rp_{y!}\) as \(t\)-exact from
affineness alone would be incorrect.
\end{perversebox}
